\subsection{Elementare Kardinalzahlgesetze und Funktionsgraphen}

\Needspace{16\baselineskip}
\FormulaThmDeltaK[Ein Funktionsgraph über einem endlichen Definitionsbereich ist endlich]{%
  \Finite D\dsep F\colon D\to A\vdash\Finite F
}{FiniteFunctionGraph}{\DeltaRow{Mengen}{A\dsep D\dsep F}}
\begin{tabproofwide}
  \proofstepwidestar[1]{\Finite D}{\rA}
  \proofstepwidestar[2]{F\colon D\to A}{\rA}
  \proofstepwidestar[2]{D\EqCard F}{\FormulaRefAuto{FunctionGraphEquinumerousDomain}[thm]{2}}
  \proofstepwidestar[1,2]{\Finite F}{\FormulaRefAuto{FiniteEqCardTransport}[thm]{1,3}}
\end{tabproofwide}

\Needspace{20\baselineskip}
\FormulaThmDeltaK[Kardinalzahl eines Funktionsgraphen mit endlichem Definitionsbereich]{%
  \Finite D\dsep F\colon D\to A\vdash\FiniteCard(F)=\FiniteCard(D)
}{FiniteFunctionGraphCardinality}{\DeltaRow{Mengen}{A\dsep D\dsep F}}
\begin{tabproofwide}
  \proofstepwidestar[1]{\Finite D}{\rA}
  \proofstepwidestar[2]{F\colon D\to A}{\rA}
  \proofstepwidestar[1,2]{\Finite F}{\FormulaRefAuto{FiniteFunctionGraph}[thm]{1,2}}
  \proofstepwidestar[2]{D\EqCard F}{\FormulaRefAuto{FunctionGraphEquinumerousDomain}[thm]{2}}
  \proofstepwidestar[1]{\NatSeg{\FiniteCard(D)}\EqCard D}{\FormulaRefAuto{FiniteCardCounts}[thm]{1}}
  \proofstepwidestar[1,2]{\NatSeg{\FiniteCard(D)}\EqCard F}{\FormulaRefAuto{EqCardTransitive}[thm]{5,4}}
  \proofstepwidestar[1]{\FiniteCard(D)\in\mathbb N}{\FormulaRefAuto{FiniteCardNatural}[thm]{1}}
  \proofstepwidestar[1,2]{\FiniteCard(F)\in\mathbb N}{\FormulaRefAuto{FiniteCardNatural}[thm]{3}}
  \proofstepwidestar[1,2]{\NatSeg{\FiniteCard(F)}\EqCard F}{\FormulaRefAuto{FiniteCardCounts}[thm]{3}}
  \proofstepwidestar[1,2]{\FiniteCard(F)=\FiniteCard(D)}{\FormulaRefAuto{FiniteCardWitnessUnique}[thm]{8,7,9,6}}
\end{tabproofwide}

\Needspace{16\baselineskip}
\noindent\begin{minipage}{\linewidth}
\FormulaThmDeltaK[Kardinalzahl der leeren Menge]{%
  \FiniteCard(\varnothing)=0
}{B27FiniteCardEmpty}{\DeltaRow{Mengen}{n}}
\end{minipage}\par
\begin{tabproofwide}
  \proofstepwidestar[]{\Finite\varnothing}{\FormulaRefAuto{FiniteEmptyAxiom}[ax]}
  \proofstepwidestar[]{\FiniteCard(\varnothing)\in\mathbb N}{\FormulaRefAuto{FiniteCardNatural}[thm]{1}}
  \proofstepwidestar[]{\NatSeg{\FiniteCard(\varnothing)}\EqCard\varnothing}{\FormulaRefAuto{FiniteCardCounts}[thm]{1}}
  \proofstepwidestar[]{0\in\mathbb N}{\FormulaRefAuto{PeanoZero}[ax]}
  \proofstepwidestar[]{\NatSeg0=\varnothing}{\FormulaRefAuto{PeanoNatSegZero}[thm]}
  \proofstepwidestar[]{\varnothing\colon\varnothing\bij\varnothing}{%
    \FormulaRefAuto{\varnothing\colon\varnothing\bij\varnothing}[thm]}
  \proofstepwidestar[]{\NatSeg0\EqCard\varnothing}{%
    \rIE{\FormulaRefAuto{a=b\vdash b=a}[thm]{5},%
      \FormulaRefAuto{A\EqCard B\coloneqq\exists F\colon A\bij B}[def]{\rEI{6}}}}
  \proofstepwidestar[]{\FiniteCard(\varnothing)=0}{%
    \FormulaRefAuto{FiniteCardWitnessUnique}[thm]{2,4,3,7}}
\end{tabproofwide}

\Needspace{16\baselineskip}
\noindent\begin{minipage}{\linewidth}
\FormulaThmDeltaK[Adjunktion erhöht die Kardinalzahl um eins]{%
  \Finite M\dsep x\notin M\vdash\FiniteCard(M\cup\{x\})=\FiniteCard(M)+1
}{B27FiniteCardAdjunction}{\DeltaRow{Mengen}{M\dsep x}}
\end{minipage}\par
\begin{tabproofwide}
  \proofstepwidestar[1]{\Finite M}{\rA}
  \proofstepwidestar[2]{x\notin M}{\rA}
  \proofstepwidestar[1]{\FiniteCard(M)\in\mathbb N}{\FormulaRefAuto{FiniteCardNatural}[thm]{1}}
  \proofstepwidestar[1]{\NatSeg{\FiniteCard(M)}\EqCard M}{\FormulaRefAuto{FiniteCardCounts}[thm]{1}}
  \proofstepwidestar[1]{\FiniteCard(M)\notin\NatSeg{\FiniteCard(M)}}{\FormulaRefAuto{PeanoNatSegEndpointNotIn}[thm]{3}}
  \proofstepwidestar[1,2]{\NatSeg{\FiniteCard(M)}\cup\{\FiniteCard(M)\}\EqCard M\cup\{x\}}{%
    \FormulaRefAuto{EqCardFreshAdjunction}[thm]{4,5,2}}
  \proofstepwidestar[1]{\NatSeg{\FiniteCard(M)+1}=\NatSeg{\FiniteCard(M)}\cup\{\FiniteCard(M)\}}{%
    \FormulaRefAuto{PeanoNatSegAddOne}[thm]{3}}
  \proofstepwidestar[1,2]{\NatSeg{\FiniteCard(M)+1}\EqCard M\cup\{x\}}{%
    \rIE{\FormulaRefAuto{a=b\vdash b=a}[thm]{7},6}}
  \proofstepwidestar[1,2]{\Finite{M\cup\{x\}}}{\FormulaRefAuto{FiniteAdjunctionAxiom}[ax]{2,1}}
  \proofstepwidestar[1,2]{\FiniteCard(M\cup\{x\})\in\mathbb N}{\FormulaRefAuto{FiniteCardNatural}[thm]{9}}
  \proofstepwidestar[1,2]{\NatSeg{\FiniteCard(M\cup\{x\})}\EqCard M\cup\{x\}}{\FormulaRefAuto{FiniteCardCounts}[thm]{9}}
  \proofstepwidestar[1]{\FiniteCard(M)+1\in\mathbb N}{\FormulaRefAuto{PeanoAddOneClosure}[thm]{3}}
  \proofstepwidestar[1,2]{\FiniteCard(M\cup\{x\})=\FiniteCard(M)+1}{%
    \FormulaRefAuto{FiniteCardWitnessUnique}[thm]{10,12,11,8}}
\end{tabproofwide}

\Needspace{24\baselineskip}
\noindent\begin{minipage}{\linewidth}
\FormulaThmDeltaK[Induktion nach der Kardinalzahl einer endlichen Menge]{%
  \begin{aligned}[t]
    &Q(0)\dsep\forall n\in\mathbb N\,(Q(n)\rightarrow Q(n+1))\dsep{}\\[-2pt]
    &\Finite M\vdash Q(\FiniteCard(M))
  \end{aligned}
}{B27FiniteCardinalityInduction}{%
  \DeltaRow{Mengen}{M\dsep n}
  \DeltaRow{Prädikat}{Q}
}
\end{minipage}\par
Die Kardinalzahl einer endlichen Menge ist eine natürliche Zahl.
Deshalb kann die durch Zahleninduktion bewiesene Aussage an dieser
Kardinalzahl ausgewertet werden.
\Needspace{12\baselineskip}
\begin{tabproofwide}
  \proofstepwidestar[1]{Q(0)}{\rA}
  \proofstepwidestar[2]{\forall n\in\mathbb N\,(Q(n)\rightarrow Q(n+1))}{\rA}
  \proofstepwidestar[3]{\Finite M}{\rA}
  \proofstepwidestar[3]{\FiniteCard(M)\in\mathbb N}{\FormulaRefAuto{FiniteCardNatural}[thm]{3}}
  \proofstepwidestar[1,2]{\forall n\in\mathbb N\,Q(n)}{\FormulaRefAuto{PeanoInductionAddOne}[thm]{1,2}}
  \proofstepwidestar[1,2,3]{Q(\FiniteCard(M))}{\rRE{4,\rUE{5}}}
\end{tabproofwide}
